% ===========================================================================
%  Chapitre 1 — Généralités sur les fonctions numériques
%  Partie I — Analyse
% ===========================================================================
\bmtchapitre{Généralités sur les fonctions numériques}{%
  Lire une image et un antécédent sur une courbe comme sur une expression,
  déterminer un ensemble de définition, établir qu'une fonction est paire,
  impaire ou périodique, et en tirer l'intervalle sur lequel il suffit de
  l'étudier.}{Analyse \textbullet\ 8 heures}

En seconde, une fonction était le plus souvent une machine à calculer : on
donnait un nombre, elle en rendait un autre. Cette année, on change de point
de vue. Avant de calculer quoi que ce soit, on commence par délimiter le
terrain — pour quels nombres la machine fonctionne-t-elle ? — puis on cherche
ce que la fonction répète, ce qu'elle reflète, ce qu'elle laisse invariant.
Une fonction paire n'a besoin d'être étudiée que sur la moitié de son
domaine ; une fonction périodique, que sur une seule période. Le chapitre
entier tient dans cette économie.

% ---------------------------------------------------------------------------
\begin{revision}[ce qui doit répondre tout de suite]
Vous aurez besoin, dès les premières pages, des acquis de seconde suivants.

\begin{enumerate}[label=\textbf{\arabic*.}, leftmargin=1.5em, itemsep=0.25em]
  \item Les six fonctions de référence : $x \mapsto ax+b$, $x \mapsto x^{2}$,
        $x \mapsto x^{3}$, $x \mapsto \dfrac{1}{x}$, $x \mapsto \sqrt{x}$ et
        $x \mapsto \abs{x}$, avec leur courbe et leur sens de variation.
  \item Résoudre une inéquation du second degré et dresser un tableau de
        signes.
  \item La condition d'existence de $\sqrt{A}$ : il faut $A \geq 0$. La
        condition d'existence de $\dfrac{1}{B}$ : il faut $B \neq 0$.
  \item Sur le cercle trigonométrique : le radian, les angles remarquables,
        et l'égalité $\cos^{2}x + \sin^{2}x = 1$.
  \item Lire les coordonnées d'un point sur un graphique et savoir ce que
        signifie « le point $A(2\,;\,5)$ appartient à la courbe ».
\end{enumerate}
\end{revision}

\begin{automatismes}
À faire de tête, sans écrire autre chose que la réponse.

\begin{multicols}{2}
\begin{enumerate}[label=\textbf{\arabic*.}, leftmargin=1.5em, itemsep=0.15em]
  \item Pour quelles valeurs de $x$ a-t-on $x - 3 \geq 0$ ?
  \item Pour quelles valeurs de $x$ a-t-on $x^{2} - 4 = 0$ ?
  \item Calculer $\abs{-7}$, puis $\abs{3 - 10}$.
  \item Calculer $\sqrt{49}$, puis $\sqrt{(-5)^{2}}$.
  \item Résoudre $2x + 1 \neq 0$.
  \item Donner le signe de $-x^{2} - 1$ selon $x$.
  \item Calculer $\cos(0)$, $\sin\!\left(\dfrac{\pi}{2}\right)$,
        $\cos(\pi)$.
  \item Développer $(-x)^{2}$, puis $(-x)^{3}$.
  \item Résoudre $x^{2} = 9$.
  \item Le nombre $-4$ a-t-il une racine carrée ?
\end{enumerate}
\end{multicols}
\end{automatismes}

% ===========================================================================
\section{Fonction, image et antécédent}

\begin{definition}[fonction numérique d'une variable réelle]
Soit $\Df$ une partie de $\R$. Définir une \textbf{fonction numérique} $f$ sur
$\Df$, c'est associer à tout réel $x$ de $\Df$ un unique réel, noté $f(x)$.
On écrit
\[
  f : \begin{array}{ccc}
        \Df & \longrightarrow & \R \\[0.15em]
        x   & \longmapsto     & f(x)
      \end{array}
\]
Le réel $f(x)$ s'appelle l'\textbf{image} de $x$ par $f$. Réciproquement, si
$y$ est un réel et si $x$ est un élément de $\Df$ tel que $f(x) = y$, on dit
que $x$ est un \textbf{antécédent} de $y$ par $f$.
\end{definition}

\begin{remarque}[le mot « unique » n'est pas décoratif]
Un réel de $\Df$ a une image et une seule. En revanche, un réel $y$ peut
avoir plusieurs antécédents, ou n'en avoir aucun : la fonction carré donne
deux antécédents à $4$, et aucun à $-1$.
\end{remarque}

\begin{visualisation}[tout se lit verticalement, puis horizontalement]
\begin{center}
\begin{tikzpicture}[scale=0.95]
  \begin{axis}[
      axis lines = middle, xlabel = {$x$}, ylabel = {$y$},
      xmin = -3.4, xmax = 3.4, ymin = -1.6, ymax = 4.6,
      xtick = {-3,-2,-1,1,2,3}, ytick = {1,2,3,4},
      width = 9.2cm, height = 6.4cm,
      tick label style = {font=\footnotesize},
      label style = {font=\footnotesize},
      axis line style = {bmtGris},
      every tick/.style = {bmtGris},
    ]
    \addplot[bmtSarcelle, very thick, samples = 200, domain = -3.1:3.1]
      {0.55*x^2 - 0.25*x};
    % image de 2
    \draw[bmtOcre, thick, dashed] (axis cs:2,0) -- (axis cs:2,1.7);
    \draw[bmtOcre, thick, dashed] (axis cs:2,1.7) -- (axis cs:0,1.7);
    \fill[bmtOcre] (axis cs:2,1.7) circle (1.9pt);
    % antécédents de 1.7
    \draw[bmtPalissandre, thick, dashed] (axis cs:-1.55,0) -- (axis cs:-1.55,1.7);
    \fill[bmtPalissandre] (axis cs:-1.55,1.7) circle (1.9pt);
    \node[font=\footnotesize, bmtOcre, anchor=west] at (axis cs:2.1,1.9) {$f(2)$};
  \end{axis}
\end{tikzpicture}
\end{center}

\vspace{0.2em}
\noindent
Pour lire l'\emph{image} de $2$ : on part de $2$ sur l'axe des abscisses, on
monte jusqu'à la courbe, on lit sur l'axe des ordonnées. Pour lire les
\emph{antécédents} de $1{,}7$ : on part de $1{,}7$ sur l'axe des ordonnées,
on se déplace horizontalement, et on relève l'abscisse de \emph{chaque} point
d'intersection rencontré. Il y en a deux ici.
\end{visualisation}

\begin{exemple}[une image, des antécédents]
Soit $f$ la fonction définie sur $\R$ par $f(x) = x^{2} - 3x$.

\textbf{Image de $4$.} On remplace : $f(4) = 4^{2} - 3 \times 4 = 16 - 12 = 4$.
L'image de $4$ est $4$.

\textbf{Antécédents de $-2$.} On résout l'équation $f(x) = -2$, c'est-à-dire
\[
  x^{2} - 3x = -2
  \quad\daccord\quad x^{2} - 3x + 2 = 0 .
\]
Le discriminant vaut $\Delta = 9 - 8 = 1 > 0$, d'où les deux solutions
$x_{1} = \dfrac{3-1}{2} = 1$ et $x_{2} = \dfrac{3+1}{2} = 2$. Le réel $-2$
admet donc exactement deux antécédents : $1$ et $2$.

\textbf{Antécédents de $-3$.} On résout $x^{2} - 3x + 3 = 0$, dont le
discriminant vaut $\Delta = 9 - 12 = -3 < 0$. Le réel $-3$ n'a aucun
antécédent par $f$.
\end{exemple}

\begin{erreurclassique}
« Calculer l'antécédent de $5$ » ne se fait pas en remplaçant $x$ par $5$ :
cela donnerait l'\emph{image} de $5$. Chercher un antécédent, c'est toujours
\emph{résoudre une équation}. Le réflexe : dès que l'énoncé dit
« antécédent », écrivez $f(x) = \ldots$ avant toute chose.
\end{erreurclassique}

\begin{entrainement}
\exo Soit $f$ définie sur $\R$ par $f(x) = 2x^{2} - 5x + 1$. Calculer
$f(0)$, $f(-1)$, $f\!\left(\dfrac{1}{2}\right)$.

\exo Soit $g$ définie sur $\R \setminus \{2\}$ par
$g(x) = \dfrac{3x+1}{x-2}$. Calculer $g(0)$, $g(3)$, $g(-1)$.

\exo Avec $f$ de l'exercice $1$ : déterminer les antécédents de $1$, puis
ceux de $-3$.

\exo Avec $g$ de l'exercice $2$ : déterminer l'antécédent éventuel de $3$.
Que constatez-vous, et comment l'expliquer ?

\exo Soit $h$ définie par $h(x) = \sqrt{x+4}$. Déterminer les antécédents
de $0$, de $3$, puis de $-1$.

\exo Une fonction $f$ vérifie $f(-2) = 7$. Traduire cette égalité par une
phrase contenant le mot « image », puis par une phrase contenant le mot
« antécédent ».
\end{entrainement}

% ===========================================================================
\section{L'ensemble de définition}

\begin{definition}[ensemble de définition]
Lorsqu'une fonction $f$ est donnée par une expression, son \textbf{ensemble de
définition}, noté $\Df_{f}$, est l'ensemble des réels $x$ pour lesquels cette
expression a un sens, c'est-à-dire pour lesquels le calcul de $f(x)$ est
possible dans $\R$.
\end{definition}

\begin{propriete}[les trois obstacles]
Dans les expressions rencontrées cette année, un réel $x$ est exclu de
l'ensemble de définition dans trois cas, et trois seulement :
\begin{enumerate}[label=\textbf{\alph*)}, leftmargin=1.8em, itemsep=0.2em]
  \item $x$ annule un \textbf{dénominateur} : le quotient $\dfrac{A}{B}$
        n'existe que si $B \neq 0$ ;
  \item $x$ rend négatif ce qui se trouve \textbf{sous une racine carrée} :
        $\sqrt{A}$ n'existe que si $A \geq 0$ ;
  \item les deux à la fois, lorsqu'une racine figure au dénominateur :
        $\dfrac{1}{\sqrt{A}}$ n'existe que si $A > 0$, l'inégalité devenant
        stricte.
\end{enumerate}
Une fonction polynôme n'est jamais concernée : son ensemble de définition
est $\R$ tout entier.
\end{propriete}

\begin{methode}{déterminer un ensemble de définition}
\begin{enumerate}[label=\textbf{\arabic*.}, leftmargin=1.6em, itemsep=0.2em]
  \item Repérer dans l'expression tous les dénominateurs et toutes les
        racines carrées. S'il n'y en a aucun, conclure : $\Df_{f} = \R$.
  \item Écrire une condition par obstacle rencontré : une équation
        $B \neq 0$, une inéquation $A \geq 0$.
  \item Résoudre chaque condition séparément. Un tableau de signes est
        souvent le chemin le plus court pour une inéquation du second degré
        ou pour un quotient.
  \item Conserver les réels qui vérifient \emph{toutes} les conditions à la
        fois : l'ensemble de définition est leur intersection.
  \item Écrire la réponse sous forme d'intervalles ou de réunion
        d'intervalles, jamais sous forme de phrase.
\end{enumerate}
\end{methode}

\begin{exemple}[un quotient]
Soit $f(x) = \dfrac{x+5}{x^{2} - x - 6}$.

Le seul obstacle est le dénominateur : il faut $x^{2} - x - 6 \neq 0$. On
résout d'abord l'équation associée : $\Delta = 1 + 24 = 25$, donc
$x_{1} = \dfrac{1-5}{2} = -2$ et $x_{2} = \dfrac{1+5}{2} = 3$.

Ces deux réels, et eux seuls, sont à exclure :
\[
  \Df_{f} = \R \setminus \{-2\,;\,3\}
          = \intervalleo{-\infty}{-2} \cup \intervalleo{-2}{3}
            \cup \intervalleo{3}{+\infty}.
\]
\end{exemple}

\begin{exemple}[une racine, puis les deux ensemble]
Soit $g(x) = \dfrac{\sqrt{x+1}}{x-4}$.

Deux obstacles cette fois. Sous la racine : $x + 1 \geq 0$, soit
$x \geq -1$. Au dénominateur : $x - 4 \neq 0$, soit $x \neq 4$.

On garde les réels vérifiant les deux conditions :
\[
  \Df_{g} = \intervallefo{-1}{4} \cup \intervalleo{4}{+\infty}.
\]
Remarquez que $-1$ est conservé — la racine y vaut $0$, ce qui est
parfaitement licite — tandis que $4$ est exclu.
\end{exemple}

\begin{erreurclassique}
Écrire $\Df_{g} = \intervallefo{-1}{+\infty}$ en « oubliant » $4$ parce que
la racine, elle, existe en $4$. Chaque obstacle doit être traité pour
lui-même : on ne conclut qu'après avoir croisé toutes les conditions.
\end{erreurclassique}

\begin{entrainement}
\exo Déterminer l'ensemble de définition de chacune des fonctions suivantes.
\begin{multicols}{2}
\begin{enumerate}[label=\textbf{\alph*)}, leftmargin=1.6em, itemsep=0.1em]
  \item $f(x) = 3x^{3} - x + 7$
  \item $f(x) = \dfrac{1}{x+5}$
  \item $f(x) = \dfrac{2x}{x^{2}-9}$
  \item $f(x) = \sqrt{2x-6}$
  \item $f(x) = \sqrt{5-x}$
  \item $f(x) = \dfrac{1}{\sqrt{x-1}}$
\end{enumerate}
\end{multicols}

\exo Même question pour
$f(x) = \dfrac{x-1}{x^{2}+1}$ et $g(x) = \dfrac{x-1}{x^{2}-1}$.
Comparer les deux résultats et dire d'où vient la différence.

\exo Déterminer l'ensemble de définition de
$f(x) = \sqrt{x^{2} - 5x + 4}$.

\exo Déterminer l'ensemble de définition de
$f(x) = \dfrac{\sqrt{3-x}}{x^{2}-4}$.

\exo Déterminer l'ensemble de définition de
$f(x) = \dfrac{1}{\abs{x} - 2}$.

\exo On considère $f(x) = \sqrt{\dfrac{x+2}{x-3}}$. Dresser le tableau de
signes du quotient, puis en déduire $\Df_{f}$.
\end{entrainement}

% ===========================================================================
\section{Sens de variation}

\begin{definition}[fonction croissante, fonction décroissante]
Soit $f$ une fonction définie sur un intervalle $I$.
\begin{itemize}[leftmargin=1.6em, itemsep=0.2em]
  \item $f$ est \textbf{croissante} sur $I$ lorsque, pour tous réels $a$ et
        $b$ de $I$ : $a \leq b \implies f(a) \leq f(b)$.
  \item $f$ est \textbf{décroissante} sur $I$ lorsque, pour tous réels $a$ et
        $b$ de $I$ : $a \leq b \implies f(a) \geq f(b)$.
\end{itemize}
Les variations sont dites \textbf{strictes} lorsque les inégalités sur les
images sont strictes. Une fonction \textbf{monotone} sur $I$ est une fonction
qui y est croissante, ou qui y est décroissante.
\end{definition}

\begin{remarque}[une variation se dit toujours \emph{sur un intervalle}]
La fonction inverse est décroissante sur $\intervalleo{-\infty}{0}$ et
décroissante sur $\intervalleo{0}{+\infty}$, mais elle n'est pas décroissante
sur leur réunion. Prenons en effet $a = -1$ et $b = 1$ : on a bien
$a \leq b$, et pourtant $f(a) = -1$ tandis que $f(b) = 1$, de sorte que
$f(a) < f(b)$ — l'inégalité n'a pas été renversée. Dire « $f$ est
décroissante sur $\R^{*}$ » est donc faux. On écrit deux lignes, pas une.
\end{remarque}

\begin{theoreme}[multiplier par une constante]
Soit $f$ une fonction monotone sur un intervalle $I$ et $\lambda$ un réel non
nul. La fonction $\lambda f$ a le même sens de variation que $f$ sur $I$ si
$\lambda > 0$, et le sens contraire si $\lambda < 0$.
\end{theoreme}

\begin{demonstration}
Supposons $f$ croissante sur $I$ et $\lambda > 0$. Soient $a$ et $b$ deux
réels de $I$ tels que $a \leq b$. Par croissance de $f$, on a
$f(a) \leq f(b)$. Multiplier les deux membres d'une inégalité par un réel
strictement positif en conserve le sens, donc
$\lambda f(a) \leq \lambda f(b)$, c'est-à-dire
$(\lambda f)(a) \leq (\lambda f)(b)$ : la fonction $\lambda f$ est croissante
sur $I$.

Si $\lambda < 0$, la même multiplication renverse l'inégalité et donne
$\lambda f(a) \geq \lambda f(b)$ : la fonction $\lambda f$ est alors
décroissante sur $I$. Le cas où $f$ est décroissante se traite mot pour mot
de la même manière.
\end{demonstration}

\begin{visualisation}[le tableau de variation résume la courbe]
\begin{center}
\begin{tikzpicture}[scale=0.95]
  \begin{axis}[
      axis lines = middle, xlabel = {$x$}, ylabel = {$y$},
      xmin = -2.6, xmax = 3.6, ymin = -3.2, ymax = 3.6,
      xtick = {-2,-1,1,2,3}, ytick = {-3,-2,-1,1,2,3},
      width = 8.6cm, height = 6.2cm,
      tick label style = {font=\footnotesize},
      label style = {font=\footnotesize},
      axis line style = {bmtGris}, every tick/.style = {bmtGris},
    ]
    \addplot[bmtSarcelle, very thick, samples = 200, domain = -2.2:3.2]
      {0.4*x^3 - 1.2*x^2 + 1};
    \fill[bmtOcre] (axis cs:0,1) circle (2.2pt);
    \fill[bmtOcre] (axis cs:2,-0.6) circle (2.2pt);
    \node[font=\scriptsize, bmtOcre, anchor=south west] at (axis cs:0.1,1.05) {$1$};
    \node[font=\scriptsize, bmtOcre, anchor=north west] at (axis cs:2.1,-0.65) {$-0{,}6$};
  \end{axis}
\end{tikzpicture}
\end{center}

\vspace{0.3em}
\noindent
La courbe monte, puis descend, puis remonte. Les deux points marqués sont
les endroits où le sens change. Le tableau de variation n'est rien d'autre
que cette lecture, écrite :

\vspace{0.4em}
\begin{center}
\renewcommand{\arraystretch}{1.15}
\begin{tabular}{|c|ccccccc|}
\hline
$x$ & $-\infty$ & & $0$ & & $2$ & & $+\infty$ \\
\hline
      & & & $1$ & & & & \\[-0.35em]
$f$   & & $\nearrow$ & & $\searrow$ & & $\nearrow$ & \\[-0.35em]
      & & & & & $-0{,}6$ & & \\
\hline
\end{tabular}
\end{center}

\vspace{0.3em}
\noindent
Les flèches se lisent de gauche à droite, et chaque flèche vaut pour un
intervalle entier : ici $\intervalleof{-\infty}{0}$, puis
$\intervalle{0}{2}$, puis $\intervallefo{2}{+\infty}$.
\end{visualisation}

\begin{entrainement}
\exo Donner, sans calcul, le sens de variation de $x \mapsto -3x + 7$ sur
$\R$, puis celui de $x \mapsto 5x^{2}$ sur $\intervallefo{0}{+\infty}$.

\exo Soit $f(x) = x^{2}$ et $g(x) = -4x^{2}$. Dresser les tableaux de
variation de $f$ et de $g$ sur $\R$ et justifier le second à l'aide du
théorème de multiplication par une constante.

\exo Soient $a$ et $b$ deux réels de $\intervallefo{0}{+\infty}$ tels que
$a \leq b$. Démontrer que $\sqrt{a} \leq \sqrt{b}$. Qu'en déduire pour la
fonction racine carrée ?

\exo Démontrer, en revenant à la définition, que la fonction
$x \mapsto \dfrac{1}{x}$ est décroissante sur $\intervalleo{0}{+\infty}$.

\exo Une fonction $f$ est croissante sur $\intervalle{-3}{5}$ et vérifie
$f(-3) = -2$ et $f(5) = 8$. Que peut-on affirmer de $f(0)$ ? Que ne peut-on
\emph{pas} affirmer ?
\end{entrainement}

% ===========================================================================
\section{Parité : quand la courbe se reflète}

Avant de parler de parité, il faut s'assurer que la question a un sens : si
$-x$ ne se trouve pas dans l'ensemble de définition alors que $x$ s'y trouve,
comparer $f(-x)$ et $f(x)$ est impossible.

\begin{definition}[ensemble symétrique par rapport à zéro]
Une partie $\Df$ de $\R$ est \textbf{symétrique par rapport à zéro} lorsque,
pour tout réel $x$, l'appartenance de $x$ à $\Df$ entraîne celle de $-x$ à
$\Df$.
\end{definition}

\begin{definition}[fonction paire, fonction impaire]
Soit $f$ une fonction dont l'ensemble de définition $\Df_{f}$ est symétrique
par rapport à zéro.
\begin{itemize}[leftmargin=1.6em, itemsep=0.2em]
  \item $f$ est \textbf{paire} lorsque $f(-x) = f(x)$ pour tout $x$ de
        $\Df_{f}$.
  \item $f$ est \textbf{impaire} lorsque $f(-x) = -f(x)$ pour tout $x$ de
        $\Df_{f}$.
\end{itemize}
\end{definition}

\begin{theoreme}[lecture géométrique de la parité]
Le plan est rapporté à un repère orthogonal $\repere$ et $\Cf_{f}$ désigne la
courbe représentative de $f$.
\begin{itemize}[leftmargin=1.6em, itemsep=0.2em]
  \item $f$ est paire \ssi\ $\Cf_{f}$ est symétrique par rapport à l'axe des
        ordonnées.
  \item $f$ est impaire \ssi\ $\Cf_{f}$ est symétrique par rapport à
        l'origine $O$ du repère.
\end{itemize}
\end{theoreme}

\begin{demonstration}[cas d'une fonction paire]
Le point $M$ de coordonnées $(x\,;\,y)$ a pour symétrique par rapport à l'axe
des ordonnées le point $M'$ de coordonnées $(-x\,;\,y)$.

Supposons $f$ paire et prenons un point $M(x\,;\,f(x))$ de $\Cf_{f}$. Son
symétrique est $M'(-x\,;\,f(x))$. Comme $\Df_{f}$ est symétrique par rapport
à zéro, le réel $-x$ appartient à $\Df_{f}$, et la parité donne
$f(-x) = f(x)$. L'ordonnée de $M'$ est donc bien $f(-x)$ : le point $M'$
appartient à $\Cf_{f}$. La courbe est donc globalement invariante par cette
symétrie.

Réciproquement, si $\Cf_{f}$ est symétrique par rapport à l'axe des
ordonnées, alors pour tout $x$ de $\Df_{f}$ le symétrique
$M'(-x\,;\,f(x))$ du point $M(x\,;\,f(x))$ appartient à $\Cf_{f}$ ; son
ordonnée vaut donc $f(-x)$, d'où $f(-x) = f(x)$.
\end{demonstration}

\begin{visualisation}[deux symétries, deux comportements]
\begin{center}
\begin{tikzpicture}[scale=0.9]
  \begin{axis}[
      name = gauche,
      axis lines = middle, xlabel = {$x$},
      xmin = -2.6, xmax = 2.6, ymin = -0.6, ymax = 4.2,
      xtick = {-2,-1,1,2}, ytick = {1,2,3,4},
      width = 6.6cm, height = 5.2cm,
      title = {\footnotesize\sffamily $f$ paire : axe $(Oy)$},
      tick label style = {font=\scriptsize},
      label style = {font=\scriptsize},
      axis line style = {bmtGris}, every tick/.style = {bmtGris},
    ]
    \addplot[bmtSarcelle, very thick, samples = 150, domain = -2.1:2.1]
      {0.85*x^2};
    \draw[bmtOcre, thick, dashed] (axis cs:-1.4,1.666) -- (axis cs:1.4,1.666);
    \fill[bmtOcre] (axis cs:1.4,1.666) circle (1.7pt);
    \fill[bmtOcre] (axis cs:-1.4,1.666) circle (1.7pt);
  \end{axis}
  \begin{axis}[
      name = droite, at = {(gauche.right of south east)}, anchor = left of south west,
      axis lines = middle, xlabel = {$x$},
      xmin = -2.6, xmax = 2.6, ymin = -3.4, ymax = 3.4,
      xtick = {-2,-1,1,2}, ytick = {-3,-2,-1,1,2,3},
      width = 6.6cm, height = 5.2cm,
      title = {\footnotesize\sffamily $f$ impaire : centre $O$},
      tick label style = {font=\scriptsize},
      label style = {font=\scriptsize},
      axis line style = {bmtGris}, every tick/.style = {bmtGris},
    ]
    \addplot[bmtSarcelle, very thick, samples = 150, domain = -2.1:2.1]
      {0.35*x^3};
    \draw[bmtOcre, thick, dashed] (axis cs:-1.6,-1.434) -- (axis cs:1.6,1.434);
    \fill[bmtOcre] (axis cs:1.6,1.434) circle (1.7pt);
    \fill[bmtOcre] (axis cs:-1.6,-1.434) circle (1.7pt);
  \end{axis}
\end{tikzpicture}
\end{center}

\vspace{0.2em}
\noindent
À gauche, les deux points marqués ont la même hauteur : on passe de l'un à
l'autre par un pliage le long de l'axe des ordonnées. À droite, le segment
qui joint les deux points marqués passe par l'origine, et celle-ci en est le
milieu : on passe de l'un à l'autre par un demi-tour autour de $O$.
\end{visualisation}

\begin{methode}{établir qu'une fonction est paire ou impaire}
\begin{enumerate}[label=\textbf{\arabic*.}, leftmargin=1.6em, itemsep=0.2em]
  \item Déterminer $\Df_{f}$ et vérifier qu'il est symétrique par rapport à
        zéro. Si ce n'est pas le cas, conclure immédiatement : la fonction
        n'est ni paire ni impaire, et il n'y a rien d'autre à écrire.
  \item Calculer $f(-x)$ en remplaçant partout $x$ par $-x$, puis simplifier
        \emph{complètement} l'expression obtenue.
  \item Comparer le résultat à $f(x)$ et à $-f(x)$ :
        \begin{itemize}[leftmargin=1.2em, itemsep=0.05em]
          \item si $f(-x) = f(x)$, la fonction est paire ;
          \item si $f(-x) = -f(x)$, la fonction est impaire ;
          \item si aucune des deux égalités n'est vraie pour tout $x$, elle
                n'est ni l'une ni l'autre — et un seul contre-exemple
                numérique suffit alors à le prouver.
        \end{itemize}
\end{enumerate}
\end{methode}

\begin{exemple}[les trois issues possibles]
\textbf{a)} $f(x) = 3x^{4} - x^{2} + 5$, définie sur $\R$, qui est symétrique
par rapport à zéro. On calcule
\[
  f(-x) = 3(-x)^{4} - (-x)^{2} + 5 = 3x^{4} - x^{2} + 5 = f(x).
\]
La fonction $f$ est paire. C'était prévisible : tous les exposants de $x$
sont pairs, et la constante $5$ se lit comme $5x^{0}$.

\textbf{b)} $g(x) = \dfrac{x}{x^{2}+1}$, définie sur $\R$. On calcule
\[
  g(-x) = \frac{-x}{(-x)^{2}+1} = \frac{-x}{x^{2}+1}
        = -\,\frac{x}{x^{2}+1} = -g(x).
\]
La fonction $g$ est impaire.

\textbf{c)} $h(x) = x^{2} + x$, définie sur $\R$. On calcule
$h(-x) = x^{2} - x$. Pour conclure, un contre-exemple suffit : $h(1) = 2$ et
$h(-1) = 0$. Comme $h(-1) \neq h(1)$ la fonction n'est pas paire, et comme
$h(-1) \neq -h(1)$ elle n'est pas impaire non plus.
\end{exemple}

\begin{propriete}[axe et centre de symétrie en un point quelconque]
Soit $f$ une fonction et $\Cf_{f}$ sa courbe dans un repère orthogonal
$\repere$. Soient $a$ et $b$ deux réels.
\begin{itemize}[leftmargin=1.6em, itemsep=0.2em]
  \item La droite d'équation $x = a$ est un \textbf{axe de symétrie} de
        $\Cf_{f}$ \ssi, pour tout réel $h$ tel que $a+h$ et $a-h$
        appartiennent à $\Df_{f}$,
        \[ f(a+h) = f(a-h). \]
  \item Le point $\Omega(a\,;\,b)$ est un \textbf{centre de symétrie} de
        $\Cf_{f}$ \ssi, pour tout réel $h$ tel que $a+h$ et $a-h$
        appartiennent à $\Df_{f}$,
        \[ f(a+h) + f(a-h) = 2b. \]
\end{itemize}
\end{propriete}

\begin{demonstration}[cas du centre de symétrie]
Le symétrique du point $M(a+h\,;\,f(a+h))$ par rapport au point
$\Omega(a\,;\,b)$ est le point $M'$ tel que $\Omega$ soit le milieu de
$[MM']$. En écrivant les coordonnées du milieu, $M'$ a pour abscisse
$2a - (a+h) = a - h$ et pour ordonnée $2b - f(a+h)$.

Dire que $\Omega$ est centre de symétrie de $\Cf_{f}$, c'est dire que $M'$
appartient à $\Cf_{f}$, c'est-à-dire que son ordonnée est l'image de son
abscisse :
\[
  2b - f(a+h) = f(a-h),
  \quad\text{soit}\quad
  f(a+h) + f(a-h) = 2b .
\]
\end{demonstration}

\begin{remarque}[la parité est un cas particulier]
En prenant $a = 0$, la première égalité devient $f(h) = f(-h)$ : c'est la
parité. En prenant $a = 0$ et $b = 0$, la seconde devient
$f(h) + f(-h) = 0$ : c'est l'imparité. Il n'y a donc qu'un seul résultat à
retenir, la parité n'en étant que la version centrée à l'origine.
\end{remarque}

\begin{entrainement}
\exo Étudier la parité de chacune des fonctions suivantes.
\begin{multicols}{2}
\begin{enumerate}[label=\textbf{\alph*)}, leftmargin=1.6em, itemsep=0.1em]
  \item $f(x) = x^{6} - 4x^{2}$
  \item $f(x) = x^{5} + 2x$
  \item $f(x) = \dfrac{1}{x^{2}+3}$
  \item $f(x) = \dfrac{x^{3}}{x^{2}-4}$
  \item $f(x) = \abs{x} - 1$
  \item $f(x) = \sqrt{x} + 1$
\end{enumerate}
\end{multicols}

\exo Pourquoi la question de la parité ne se pose-t-elle pas pour
$f(x) = \dfrac{1}{x-2}$ ? Répondre en une phrase, en citant la condition qui
n'est pas remplie.

\exo Soit $f(x) = x^{2} - 6x + 5$. Démontrer que la droite d'équation
$x = 3$ est un axe de symétrie de la courbe de $f$.

\exo Soit $g(x) = \dfrac{2x-1}{x-1}$, définie sur $\R \setminus \{1\}$.
Démontrer que le point $\Omega(1\,;\,2)$ est un centre de symétrie de la
courbe de $g$.

\exo Soit $f$ une fonction paire et $g$ une fonction impaire, toutes deux
définies sur $\R$. Que peut-on dire de la parité de $f + g$ ? de $f \times g$ ?
Justifier.

\exo Démontrer qu'une fonction définie sur $\R$, à la fois paire et impaire,
est nécessairement la fonction nulle.
\end{entrainement}

% ===========================================================================
\section{Périodicité : quand la courbe se répète}

\begin{definition}[fonction périodique]
Soit $f$ une fonction définie sur $\Df_{f}$ et $T$ un réel strictement
positif. La fonction $f$ est \textbf{périodique de période $T$} lorsque, pour
tout $x$ de $\Df_{f}$, les réels $x+T$ et $x-T$ appartiennent à $\Df_{f}$ et
\[
  f(x+T) = f(x).
\]
Le plus petit réel strictement positif vérifiant cette égalité, lorsqu'il
existe, s'appelle la \textbf{période} de $f$.
\end{definition}

\begin{propriete}[répétition et réduction de l'étude]
Si $f$ est périodique de période $T$, alors pour tout entier naturel $k$ et
tout $x$ de $\Df_{f}$ :
\[
  f(x + kT) = f(x).
\]
Il suffit donc d'étudier $f$ sur un intervalle de longueur $T$ — par exemple
$\intervallefo{0}{T}$ — et de reproduire la portion de courbe obtenue par
translations successives de vecteur $T\veci$ pour obtenir la courbe entière.
\end{propriete}

\begin{demonstration}
L'égalité est vraie pour $k = 0$ puisqu'elle s'écrit alors $f(x) = f(x)$.
Elle est vraie pour $k = 1$ par définition de la périodicité. Pour $k = 2$,
on applique deux fois la définition :
\[
  f(x + 2T) = f\big((x+T) + T\big) = f(x+T) = f(x),
\]
et l'on voit que le procédé se poursuit de proche en proche : chaque pas
ajoute un $T$ et redonne la même valeur.
\end{demonstration}

\begin{exemple}[les fonctions du cercle trigonométrique]
Les fonctions $x \mapsto \cos x$ et $x \mapsto \sin x$ sont définies sur $\R$
et périodiques de période $2\pi$ : un tour complet du cercle
trigonométrique ramène au même point, donc aux mêmes coordonnées.

La fonction $f(x) = \cos(3x)$ est également périodique. Cherchons sa période.
Pour tout réel $x$,
\[
  f\!\left(x + \frac{2\pi}{3}\right)
  = \cos\!\left(3x + 2\pi\right) = \cos(3x) = f(x),
\]
donc $\dfrac{2\pi}{3}$ est une période de $f$.
\end{exemple}

\begin{visualisation}[un motif, puis des copies]
\begin{center}
\begin{tikzpicture}
  \begin{axis}[
      axis lines = middle, xlabel = {$x$},
      xmin = -7.2, xmax = 7.2, ymin = -1.8, ymax = 1.8,
      xtick = {-6.2832, -3.1416, 3.1416, 6.2832},
      xticklabels = {$-2\pi$, $-\pi$, $\pi$, $2\pi$},
      ytick = {-1,1},
      width = 12cm, height = 4.6cm,
      tick label style = {font=\footnotesize},
      label style = {font=\footnotesize},
      axis line style = {bmtGris}, every tick/.style = {bmtGris},
    ]
    \addplot[bmtGrisClair, very thick, samples = 250, domain = -7:7]
      {sin(deg(x))};
    \addplot[bmtSarcelle, very thick, samples = 150, domain = 0:6.2832]
      {sin(deg(x))};
    \draw[bmtOcre, thick, dashed] (axis cs:0,-1.5) -- (axis cs:0,1.5);
    \draw[bmtOcre, thick, dashed] (axis cs:6.2832,-1.5) -- (axis cs:6.2832,1.5);
  \end{axis}
\end{tikzpicture}
\end{center}

\vspace{0.2em}
\noindent
La portion foncée, tracée sur $\intervallefo{0}{2\pi}$, contient toute
l'information. Le reste de la courbe, en clair, s'en déduit par translations
de vecteur $2\pi\veci$ vers la droite et vers la gauche. C'est là tout
l'intérêt de repérer une période : l'étude se fait une fois, le tracé se
répète.
\end{visualisation}

\begin{methode}{réduire l'intervalle d'étude}
Lorsqu'une fonction est à la fois périodique et paire (ou impaire), les deux
propriétés se cumulent et l'intervalle d'étude se réduit deux fois.
\begin{enumerate}[label=\textbf{\arabic*.}, leftmargin=1.6em, itemsep=0.2em]
  \item La périodicité de période $T$ ramène l'étude à un intervalle de
        longueur $T$, centré si possible en zéro :
        $\intervalle{-\frac{T}{2}}{\frac{T}{2}}$.
  \item La parité, ou l'imparité, permet alors de n'étudier que la moitié
        droite : $\intervalle{0}{\frac{T}{2}}$.
  \item On trace la courbe sur ce dernier intervalle, on complète par la
        symétrie donnée par la parité, puis on reproduit le motif obtenu par
        translations de vecteur $T\veci$.
\end{enumerate}
\end{methode}

\begin{exemple}[cumuler les deux réductions]
Soit $f(x) = \cos(2x)$, définie sur $\R$.

\textbf{Périodicité.} Pour tout réel $x$, $f(x + \pi) = \cos(2x + 2\pi)
= \cos(2x) = f(x)$. La fonction est périodique de période $\pi$, et l'on
peut se limiter à $\intervalle{-\frac{\pi}{2}}{\frac{\pi}{2}}$.

\textbf{Parité.} Pour tout réel $x$, $f(-x) = \cos(-2x) = \cos(2x) = f(x)$,
car la fonction cosinus est paire. La fonction $f$ est donc paire, et l'on
peut se limiter à $\intervalle{0}{\frac{\pi}{2}}$.

D'un intervalle infini, on est passé à un intervalle de longueur
$\dfrac{\pi}{2}$.
\end{exemple}

\begin{entrainement}
\exo Démontrer que la fonction $x \mapsto \sin(4x)$ est périodique et
déterminer sa période.

\exo Même question pour $x \mapsto \cos\!\left(\dfrac{x}{2}\right)$.

\exo Soit $f(x) = \sin x + \cos x$. Démontrer que $f$ est périodique de
période $2\pi$. Cette fonction est-elle paire ? impaire ?

\exo Soit $f(x) = \cos^{2} x$. Démontrer que $\pi$ est une période de $f$,
puis étudier sa parité et en déduire l'intervalle d'étude le plus court.

\exo Une fonction $f$ est périodique de période $3$ et vérifie $f(1) = 5$.
Calculer $f(4)$, $f(7)$ et $f(-2)$.

\exo Soit $f$ définie sur $\R$, paire et périodique de période $T$.
Démontrer que la droite d'équation $x = \dfrac{T}{2}$ est un axe de symétrie
de la courbe de $f$.
\end{entrainement}

% ===========================================================================
% BANQUE D'EXERCICES DE FIN DE CHAPITRE
% À compléter après le dépouillement des sujets d'examen de fin de première
% (probatoire du Cameroun séries C et D, compositions d'Afrique francophone).
% Blocs attendus : \begin{annales} ... \end{annales} puis \begin{plusloin}.
% ===========================================================================

\begin{jecherche}[une courbe reconstituée]
On considère une fonction $f$ définie sur $\R$, impaire et périodique de
période $4$. On sait de plus que, pour tout $x$ de $\intervalle{0}{2}$,
$f(x) = x(2-x)$.

\begin{enumerate}[label=\textbf{\arabic*.}, leftmargin=1.6em, itemsep=0.2em]
  \item Calculer $f(0)$, $f(1)$ et $f(2)$.
  \item Déterminer $f(-1)$, puis $f(3)$, puis $f(5)$, en précisant à chaque
        fois laquelle des deux hypothèses est utilisée.
  \item Tracer la courbe de $f$ sur $\intervalle{0}{2}$, puis la compléter
        sur $\intervalle{-4}{6}$.
  \item La droite d'équation $x = 1$ est-elle un axe de symétrie de cette
        courbe ? Justifier.
\end{enumerate}
\end{jecherche}

\begin{jemeteste}
\begin{enumerate}[label=\textbf{\arabic*.}, leftmargin=1.6em, itemsep=0.25em]
  \item Déterminer l'ensemble de définition de
        $f(x) = \dfrac{\sqrt{x+3}}{x^{2}-16}$.
  \item Soit $g(x) = 2x^{2} - 8x + 1$. Déterminer les antécédents de $1$ par
        $g$.
  \item Étudier la parité de $h(x) = \dfrac{x^{4}+1}{x^{2}}$.
  \item Démontrer que la droite d'équation $x = 2$ est un axe de symétrie de
        la courbe de $g$ définie à la question $2$.
  \item Soit $u(x) = \sin(3x)$. Déterminer une période de $u$, puis étudier
        sa parité, et en déduire le plus petit intervalle sur lequel il
        suffit d'étudier $u$.
\end{enumerate}
\end{jemeteste}

\begin{notepourlaclasse}
Les huit heures se répartissent naturellement en quatre séances de deux
heures : image et antécédent avec l'ensemble de définition, puis le sens de
variation, puis la parité, puis la périodicité avec la réduction de
l'intervalle d'étude.

Le point qui coûte le plus de temps chaque année n'est pas la parité mais la
\emph{vérification préalable de la symétrie de l'ensemble de définition}. Les
élèves calculent $f(-x)$ sans y penser et concluent sur une fonction dont le
domaine ne s'y prête pas. L'exercice 2 de la banque sur la parité est fait
pour cela : le faire traiter au tableau, pas en devoir.

Dans une classe à grand effectif, la démonstration du centre de symétrie peut
être menée collectivement sur une figure au tableau, le calcul des
coordonnées du symétrique étant la seule étape réellement nouvelle ; le cas
de l'axe de symétrie est alors laissé en exercice.
\end{notepourlaclasse}
